\section{Related Work}
\label{ch:related}

The work in this thesis sits at the intersection of three lines of research: controllable text generation, energy-based and sampling-based decoding, and the study of the geometry and behaviour of language-model representations. This section reviews each in turn. Rather than cataloguing methods, it tells the story of how the field arrived at energy-based sampling as an approach to control, because understanding that trajectory is what makes the unexamined assumption at its heart visible, and it is that assumption which the thesis tests. The section closes by stating precisely the gap in the literature that the work addresses.

\subsection{Controllable Text Generation}
\label{sec:rel-ctg}

Approaches to controllable generation can be organized by how much they change the underlying model, and reviewing them in that order reveals a progression toward ever lighter-touch intervention, of which energy-based sampling is the logical endpoint.

At one extreme are the \textbf{fine-tuning methods}, which adjust the model's parameters so that the desired behaviour is internalized. The earliest of these simply continued training on attribute-specific data, and the modern and dominant form is reinforcement learning from human feedback \citep{ouyang2022training}, which optimizes the model against a learned reward that encodes human preferences. These methods work, and the widely used aligned assistants are their fruit, but the control they produce is fixed once training is complete. A new attribute requires new data and a new training run, and the cost of that adaptation is precisely what motivates the search for methods that impose control at inference time. Conditioning methods such as CTRL \citep{keskar2019ctrl} occupy a middle ground, prepending control codes to the input so that a single model can be steered among a fixed set of attributes anticipated during training, but they too cannot accommodate a genuinely novel constraint after the fact.

At the other extreme are the \textbf{decoding-time methods}, which leave the model's parameters entirely untouched and intervene only during generation. These are the methods most relevant to this thesis, because they are the ones that promise the inference-time flexibility that fine-tuning cannot offer. Plug and Play Language Models \citep{dathathri2020plug} perturb the hidden states of a frozen model at each step in the direction indicated by an attribute classifier's gradient, nudging the generation toward the desired attribute without altering any weight. FUDGE \citep{yang2021fudge} and GeDi \citep{krause2021gedi} take a complementary approach, reweighting the next-token distribution using a discriminator that predicts, from a prefix, whether the eventual completion will bear the desired attribute. These methods are ingenious and effective within their scope, but they share a structural feature that limits them: they remain autoregressive. They modify the local, per-token decision, but they retain the strictly left-to-right generation order and therefore inherit its fundamental inability to revise an earlier token in light of a later one. A method that could edit any part of the sequence in light of the whole would be strictly more powerful, and it is the pursuit of that capability that leads to the sampling-based methods.

\subsection{Energy-Based and Sampling-Based Decoding}
\label{sec:rel-energy}

The decisive conceptual move, and the one that defines the family of methods this thesis studies, is to abandon left-to-right generation entirely and to cast decoding as sampling from an energy defined over whole sequences. Once a sequence is treated as a single object drawn from a distribution, rather than as a stream of tokens produced in order, a global constraint can be imposed on the completed object directly, and any position in the sequence can be revised in light of any other. This is the promise that makes the approach attractive, and several influential systems realize it.

COLD decoding \citep{qin2022cold} performs energy-based constrained generation by running Langevin dynamics on a continuous relaxation of the sequence, combining a fluency energy derived from a language model with task-specific constraint energies, and it demonstrated the approach on tasks such as constrained generation and abductive reasoning. MuCoLa \citep{kumar2022gradient} samples in the continuous embedding space and projects back to tokens, and it is the closest published antecedent of the continuous sampler studied in this thesis; its formulation of the sampling problem, and its use of a projection to return to discrete text, are exactly the ingredients whose geometric consequences Section~\ref{ch:results} examines. Related efforts include gradient-based inference for lexically constrained generation and infilling \citep{sha2020gradient}, and a body of work on discrete samplers adapted from the energy-based-model literature, including the gradient-informed discrete proposals of \citet{grathwohl2021oops} and the discrete Langevin proposals of \citet{zhang2022langevin} on which the discrete sampler of this thesis is directly based.

Reading across this literature, a feature recurs that is rarely brought to the foreground, and it is the feature this thesis makes central. The Metropolis--Hastings correction, which as Section~\ref{sec:bg-mh} explained is what separates genuine sampling from biased optimization, is frequently omitted. Methods are described in the language of sampling from an energy while in practice running a noise-augmented gradient descent, and where the correction is included, the difficulty of computing its reverse-proposal term on a landscape defined by projection is seldom examined. This omission is understandable, because including the correction faithfully often makes these methods stop working, but the reason it makes them stop working is itself informative, and recovering that reason is one of the contributions of this thesis. Taking the correction seriously, implementing it faithfully for both a discrete and a continuous sampler, and measuring exactly what its inclusion reveals is a lens through which the failure of the whole approach comes into focus.

A distinct branch of the same literature deserves to be isolated here, because it anticipates the diagnosis this thesis reaches. Not every method that samples from a frozen-language-model energy uses the \emph{gradient} of that energy. Mix-and-Match \citep{mireshghallah2022mixmatch} performs controllable generation by Metropolis--Hastings over a product-of-experts energy, proposing token substitutions from an auxiliary masked language model and scoring them with the frozen energy, so that the energy is only ever evaluated, never differentiated. The twisted sequential Monte Carlo methods of \citet{lew2023sequential} and \citet{zhao2024probabilistic} likewise sample from an intractable posterior defined by a frozen model using importance weights and resampling rather than energy gradients. That these gradient-free samplers operate successfully on essentially the same frozen-LM energies that defeat the gradient-guided methods is, in retrospect, a strong hint about where the difficulty lies, and this thesis makes the hint concrete: a gradient-free Gibbs sampler run on exactly the energy the Langevin samplers navigate, reported in Section~\ref{sec:results-baselines}, recovers coherent text on the same benchmark. That in-platform result is the clean separation between evaluating the energy, which works, and differentiating it, which does not, and it is what licenses narrowing the thesis's claim to the gradient rather than the energy.

\subsection{Amortized Sampling and GFlowNets}
\label{sec:rel-gflownet}

Generative Flow Networks were introduced by \citet{bengio2021flow} as a method for sampling discrete, compositional objects in proportion to a reward, with the explicit motivation of producing diverse high-reward samples rather than the single mode that reward-maximizing reinforcement learning tends to collapse onto. The framework was subsequently formalized and given a theoretical foundation \citep{bengio2023gflownet}, and a family of training objectives was developed to address the practical difficulty of assigning credit along long generative trajectories, including trajectory balance \citep{malkin2022trajectory} and its partial-episode sub-trajectory refinement \citep{madan2023learning}. The application most directly relevant here is that of \citet{hu2024amortizing}, who framed a range of language tasks, including infilling and reasoning, as sampling from an intractable posterior defined by a frozen language model, and trained a GFlowNet to perform that sampling, reporting that the resulting policies yield diverse samples that transfer better to downstream use than those obtained by reward maximization.

The GFlowNet component of this thesis builds directly on that formulation and on the accompanying codebase, using the same modified sub-trajectory objective and the same parameter-efficient adaptation of a frozen base model \citep{hu2022lora}. The question posed here, however, is narrower and more critical than the one those works set out to answer. Their concern is whether amortized sampling is useful, and they show that it is, on tasks where the reward defines a well-behaved target. The concern here is whether amortized sampling \emph{escapes a specific pathology}, namely that of the frozen autoregressive likelihood used as an energy over free-form sequences, which the first half of this thesis shows defeats the gradient-guided samplers. These questions are compatible, and their answers turn out to be compatible too: the difference lies entirely in whether the reward defines a landscape with a coherent-text optimum, and the contribution of Section~\ref{ch:results} is to show, by direct experiment, that the frozen likelihood does not, so that amortization inherits rather than escapes the problem.

\subsection{The Geometry and Behaviour of Language-Model Representations}
\label{sec:rel-geometry}

A final and complementary line of work concerns the representations in which these samplers operate and the behaviour of the likelihood they optimize, and it supplies part of the explanation for the results. Two phenomena from this literature are directly relevant.

The first is \textbf{anisotropy}. A well-documented property of the embeddings produced by Transformer language models, both contextual and static, is that they do not spread uniformly through the representation space but instead occupy a narrow cone, so that two randomly chosen token embeddings have high average cosine similarity rather than being roughly orthogonal as one might naively expect \citep{gao2019representation, ethayarajh2019contextual}. This has a direct consequence for any method that treats Euclidean or cosine distance in embedding space as a meaningful measure of similarity, because anisotropy compresses the dynamic range of those distances: the difference between a good and a bad token, measured in the embedding metric, is smaller than the geometry would suggest. This thesis measures the anisotropy of both models it studies and shows that it renders Euclidean distance an unreliable evaluation metric, which is one reason the experiments adopt a divergence-based measure of contextual fit instead, and it shows in addition that the anisotropy scale differs enough between models to make a single step-size schedule non-transferable, a consequence the prior work did not draw.

The second is the failure of high-probability text to be high-quality text. \citet{holtzman2020curious} demonstrated that the most probable sequences under a neural language model are frequently degenerate, repetitive, and unnatural, and that maximization-based decoding such as beam search produces markedly worse text than sampling from the same model. This observation is central to the interpretation of the present results, because it means that any procedure which succeeds in concentrating probability mass on the lowest-energy regions of the frozen-likelihood energy is thereby concentrating it on degenerate text. This thesis reproduces the phenomenon quantitatively on its own models, terms it the \emph{likelihood trap} in the context of energy-based sampling, and uses it to argue that the objective the samplers pursue is itself undesirable, quite apart from whether they can pursue it effectively.

\subsection{Diffusion Models for Text and the Score Function}
\label{sec:rel-diffusion}

A final body of work is directly relevant not to the methods this thesis evaluates but to the positive control it proposes, and so it is reviewed here to prepare the discussion of Section~\ref{ch:discussion}. Diffusion models, which generate by gradually reversing a noising process, have become the dominant approach to image generation and have been adapted to text in two broad families. The first embeds tokens into a continuous space and runs a continuous diffusion there, as in Diffusion-LM \citep{li2022diffusion}, which demonstrated controllable generation by applying classifier guidance to the continuous latents at each denoising step. The second operates directly on discrete tokens, as in the score-entropy discrete diffusion of \citet{lou2024discrete}, which learns the ratios of the data distribution over a discrete state space and has narrowed much of the quality gap between diffusion and autoregressive language models.

What makes these models pertinent to this thesis is not their generation quality but their training objective. The foundational result connecting them to the samplers studied here is that of \citet{vincent2011connection}: learning to denoise a corrupted input is mathematically equivalent to learning the score of the data distribution, the gradient of its log-density, and it is this equivalence that gives score-based generative modeling its name and its mechanism \citep{song2019generative, ho2020denoising}. A diffusion model, unlike an autoregressive one, is therefore trained precisely to provide the gradient field that Langevin dynamics requires, at every noise level and across the whole space rather than only on the data manifold. This is exactly the property that the results of this thesis find the autoregressive likelihood to lack, which is why, as Section~\ref{ch:discussion} develops, a diffusion language model is the natural instrument for a positive control: it holds the sampling method fixed and changes only the training objective that shapes the energy, and so it can isolate that objective as the cause of the failure this thesis diagnoses.

\subsection{The Gap Addressed by This Thesis}
\label{sec:rel-gap}

The literature reviewed above establishes energy-based sampling as an attractive and actively pursued route to controllable generation, and it supplies both the samplers this thesis studies and the amortized alternative it examines. What the literature does not supply is a direct, controlled test of the assumption on which the entire enterprise rests, namely that the gradient of a frozen autoregressive likelihood is a usable search direction on the discrete text manifold, together with a mechanistic account of what happens when that assumption fails. Negative observations, where they appear, tend to be reported in passing, as brittleness, as sensitivity to hyperparameters, or as a need for careful tuning, rather than isolated as a phenomenon and explained. This thesis addresses that gap directly. It designs an ablation that cleanly separates the direction of the gradient from its magnitude, and thereby tests the assumption rather than a proxy for it; it constructs diagnostic experiments that identify the geometric, statistical, and analytic mechanisms behind the result; and it shows that the same mechanism survives amortization, which locates the cause in the training objective of the language model rather than in any particular sampling algorithm. The result is not merely the observation that these methods fail, which is implicit in the difficulty the field has had with them, but an explanation of why, supported at every step by measurement.
